
\chapter{攻读硕士学位期间取得的研究成果}

{
\zihao{5}
% \setstretch{1.65}
\vspace{2pt}
一、已发表（包括已接受待发表）的论文，以及已投稿、或已成文打算投稿、或拟成文投稿的论文情况\underline{\textbf{（只填写与学位论文内容相关的部分）：}}
\vspace{1em}
% ==================== 表格区域 ====================
% 定义列格式：m代表垂直居中，后面的长度是列宽，<{\centering}代表水平居中
\begin{table}[h]
    \centering
    \zihao{-5}
    \setlength{\tabcolsep}{3pt} % 减小列间距
    \begin{tabular}{|m{0.8cm}<{\centering}|m{3.2cm}<{\centering}|m{3.5cm}<{\centering}|m{3.2cm}<{\centering}|m{1.8cm}<{\centering}|m{1.8cm}<{\centering}|m{1.5cm}<{\centering}|}
        \hline
        序号 & 
        作者（全体作者，按顺序排列） & 
        题目 & 
        发表或投稿刊物名称、级别 & 
        发表的卷期、年月、页码 & 
        与学位论文哪一部分（章、节）相关 & 
        被索引收录情况 \\ 
        \hline
        
        % 第一行数据
        1 & 
        \textbf{Dong Liu}, Yifan Yang, Zixiong Huang, Yuxin Gao, Mingkui Tan & 
        CHRIS: Clothed Human Reconstruction with Side View Consistency & 
        IEEE International Conference on Multimedia \& Expo \par \vspace{0.5em} \textbf{CCF-B 类会议} & 
        2025 & 
        第四章 & 
        已接收 \\ 
        \hline
        
        % 第二行数据
        2 & 
        Yifan Yang, \textbf{Dong Liu}, Shuhai Zhang, Zeshui Deng, Zixiong Huang, Mingkui Tan & 
        HiLo: Detailed and Robust 3D Clothed Human Reconstruction with High-and Low-Frequency Information of Parametric Models & 
        The IEEE / CVF Computer Vision and Pattern Recognition Conference \par \vspace{0.5em} \textbf{CCF-A 类会议} & 
        2024 & 
        第二章 & 
        EI 收录 \\ 
        \hline
    \end{tabular}
\end{table}

\vspace{0.5em}

% ==================== 注释区域 ====================
\noindent 注：在“发表的卷期、年月、页码”栏：

\noindent 1. 如果论文已发表，请填写发表的卷期、年月、页码；

\noindent 2. 如果论文已被接受，填写将要发表的卷期、年月；

\noindent 3. 以上都不是，请据实填写“已投稿”，“拟投稿”。

\noindent 不够请另加页。

\vspace{1.5em}

% ==================== 第二部分 ====================
\noindent 二、与学位内容相关的其它成果（包括专利、著作、获奖项目等）

\vspace{0.5em}

\noindent 专利成果：

\begin{itemize}
    \item 杜卿、\textbf{刘东}、杨逸凡、谭明奎. 基于高频局部特征的三维穿衣人体重建方法、装置及介质. 
    CN202311764613.8.
\end{itemize}
}